
\documentclass{article}
\usepackage{colortbl}
\usepackage{makecell}
\usepackage{multirow}
\usepackage{supertabular}

\begin{document}

\newcounter{utterance}

\centering \large Interaction Transcript for game `guesswhat', experiment `Level\_3', episode 2 with reka{-}flash{-}3.1{-}t0.0.
\vspace{24pt}

{ \footnotesize  \setcounter{utterance}{1}
\setlength{\tabcolsep}{0pt}
\begin{supertabular}{c@{$\;$}|p{.15\linewidth}@{}p{.15\linewidth}p{.15\linewidth}p{.15\linewidth}p{.15\linewidth}p{.15\linewidth}}
    \# & $\;$A & \multicolumn{4}{c}{Game Master} & $\;\:$B\\
    \hline

    \theutterance \stepcounter{utterance}  
    & & \multicolumn{4}{p{0.6\linewidth}}{
        \cellcolor[rgb]{0.9,0.9,0.9}{
            \makecell[{{p{\linewidth}}}]{
                \texttt{\tiny{[P1$\langle$GM]}}
                \texttt{Let's play a guessing game. Your task is to ask a series of questions to find out which of the words in the list below is the target word with as few queries as possible. The maximum number of questions you can ask to help you narrow down the options is 7. Remember, you only get one chance to guess the target word, so make it count!} \\
\\ 
\texttt{The rules of the game are: } \\
\\ 
\texttt{1. Only ask one question per turn. DO NOT ASK MORE THAN ONE QUESTION.} \\
\\ 
\texttt{2. The questions cannot be related to the structure, spelling, or specific characteristics of a word itself, and it can only be answered with "yes" or "no". } \\
\\ 
\texttt{3. Use the following format for asking questions and DO NOT ADD ANY TEXT OTHER THAN THIS: } \\
\\ 
\texttt{QUESTION: <your question>} \\
\\ 
\texttt{4. You have only one opportunity to guess the target word. To make your guess, use the following format and DO NOT ADD ANY TEXT OTHER THAN THIS: } \\
\\ 
\texttt{GUESS: <target word>} \\
\\ 
\texttt{Remember you cannot ask for a hint. Make your question brief. Let’s start!} \\
\\ 
\texttt{List of candidate words: {[}'Nephew', 'Uncle', 'Grandfather', 'Grandmother', 'Father', 'Mother', 'Wife', 'Husband'{]}} \\
            }
        }
    }
    & \\ \\

    \theutterance \stepcounter{utterance}  
    & \multicolumn{4}{p{0.6\linewidth}}{
        \cellcolor[rgb]{0.9,0.9,0.9}{
            \makecell[{{p{\linewidth}}}]{
                \texttt{\tiny{[P1$\rangle$GM]}}
                \texttt{<reasoning>} \\
\texttt{Let's see. I need to figure out which word from the list is the target using up to 7 yes/no questions. The list includes family relations: Nephew, Uncle, Grandfather, Grandmother, Father, Mother, Wife, Husband.} \\
\\ 
\texttt{First, I should categorize these words to see if there's a way to split them into groups. Let's think about gender. Some are male, some female. For example, Father, Uncle, Grandfather, Husband are male. Mother, Grandmother, Wife are female. Nephew is male, but he's a younger son's son. So maybe first question: Is the target a male? That could split the list into two groups. Let's count: male options are Father, Uncle, Grandfather, Husband, Nephew (5). Female are Mother, Grandmother, Wife (3). If the answer is yes, then narrow down to 5; if no, 3. But maybe there's a better way.} \\
\\ 
\texttt{Alternatively, think about generational levels. Grandfather and Grandmother are grandparents, Uncle and Aunt (but Aunt isn't in the list) are siblings of parents. Father and Mother are parents, Husband and Wife are spouses, Nephew is a child of a sibling. So maybe first question: Is the target a grandparent? If yes, then Grandfather or Grandmother. If no, then the rest. Let's see. If the answer is yes, then two options. If no, then 6 options left. But maybe that's too narrow. Alternatively, maybe ask if it's a parent. Parents are Father and Mother. If yes, then two options. If no, then the rest. Let's see. Let's try that.} \\
\\ 
\texttt{Question 1: Is the target a parent? Possible answers: Father or Mother if yes. If no, then the other six. Let's say the answer is no. Then we have to choose from Uncle, Grandfather, Grandmother, Husband, Wife, Nephew. Next question: Is the target a grandparent? If yes, then Grandfather or Grandmother. If no, then the remaining four. If answer is yes, then two options. If no, then four. Let's say answer is no. Then next question: Is the target a spouse? Spouses are Husband and Wife. If yes, then two options. If no, then Uncle, Neph} \\
            }
        }
    }
    & & \\ \\

    \theutterance \stepcounter{utterance}  
    & & & \multicolumn{2}{p{0.3\linewidth}}{
        \cellcolor[rgb]{0.9,0.9,0.9}{
            \makecell[{{p{\linewidth}}}]{
                \texttt{\tiny{[GM$|$GM]}}
                \texttt{Invalid format. Guesser must use the form 'QUESTION: ' or 'GUESS: '.} \\
            }
        }
    }
    & & \\ \\

    \theutterance \stepcounter{utterance}  
    & & & \multicolumn{2}{p{0.3\linewidth}}{
        \cellcolor[rgb]{0.9,0.9,0.9}{
            \makecell[{{p{\linewidth}}}]{
                \texttt{\tiny{[GM$|$GM]}}
                \texttt{abort game} \\
            }
        }
    }
    & & \\ \\

\end{supertabular}
}

\end{document}
